\section{Reservation and three physical programs}\label{simp:score-primitives}

We have $n$ normalized nonnegative monotone submodular valuations
$f_i:2^M\to\mathbb Q_{\ge0}$. Reserve a maximum-product positive singleton
matching $i\mapsto h_i$ and put $H_0=\{h_i\}$, $M'=M\setminus H_0$,
$d=|M'|$. If such a matching does not exist, $\OPT=0$: submodularity
implies $f_i(S)\le\sum_{j\in S}f_i(j)$, so positive disjoint bundles
would supply positive singleton representatives. The case $|M|<n$ is
also immediate. Rescale each active row by $f_i(M')=1$; rows with
$f_i(M')=0$ need no residual update. Thus $1/d\le B_i\le1$ on active
rows, where $B_i=\max_{g\in M'}f_i(g)$.

Put $p_{ij}=f_i(h_j)$ and $p_i^0=p_{ii}>0$. Give arc $i\to j$ weight
$p_{ij}/p_i^0$ and let $s_j$ be the largest product of a simple path
ending at $j$, or one. Every cycle has product at most one, or it would
improve the reservation. Multiplicative Bellman--Ford computes $s_j$
in $n-1$ rounds. Set $\alpha_i=p_i^0/s_i$ and $e^{\rho_j}=s_j$.
Path extension, and a path followed by a best residual singleton, give
\begin{equation}\label{na:reservation-dual}
 B_i\le\alpha_i,\qquad p_{ij}\le\alpha_i e^{\rho_j},\qquad
 p_{ii}=\alpha_i e^{\rho_i},\qquad\rho_j\ge0.
\end{equation}
These protection data are rational; the algorithm need not compute $\rho$.
For independent $R_i\sim R(x_i)$, write
$F_i=\mathbb Ef_i(R_i)$ and $d_{ij}=\mathbb Ef_i(R_i+h_j)-p_{ij}$.
Then $0\le d_{ij}\le F_i\le p_{ij}+d_{ij}$, and every residual
marginal of the raw valuation or a contraction is at most $\alpha_i$.

\subsection{One fixed mixture and its private profile}

Define
\[
 h^{C}_{p,\alpha}(X,V)=\max\{p,\min(V,X/\theta),\theta\alpha\},\quad
 h^{A}_{p,\alpha}(X,V)=\max\{X,p,aV,\theta\alpha\},\quad
 \mathcal H_{p,\alpha}=\prod_c h_{c,p,\alpha}^{\omega_c}.
\]
All entries below are exact rationals; the weights sum to one.
\begin{center}
\begin{tabular}{crrr}
\toprule Type & $a$ & $\theta$ & $\omega$\\\midrule
$C$ & --- & $1$ & $.166466453839$\\
$A$ & $.75$ & $.592592592592$ & $.048883852144$\\
$A$ & $.80$ & $.488281250000$ & $.044226298717$\\
$A$ & $.90$ & $.258117479171$ & $.032834490836$\\
$A$ & $.925$ & $.196045038119$ & $.021914693827$\\
$A$ & $.999$ & $.002716922574$ & $.254146426697$\\
$C$ & --- & $.0005$ & $.156660002590$\\
$C$ & --- & $.5$ & $.274867781350$\\\bottomrule
\end{tabular}
\end{center}
Let $\theta_0=1/2000$, $\Gamma=\prod_c\theta_c^{\omega_c}$,
$H(s)=\prod_c\max(s,\theta_c)^{\omega_c}$, and $f=\log H$. The private score is
\begin{equation}\label{na:private}
 T_{i0}=\alpha_iG(F_i/\alpha_i),\qquad
 \log G(m)=\min\{\mathbb Ef(Y):Y\ge0,\ \mathbb EY=m,
                  |\operatorname{supp}Y|\le2,\ \operatorname{width}(Y)\le1\}.
\end{equation}
It has $G(0)=\Gamma>.0487864347$ and $G(m)\ge c_0m$, $c_0=.3999999$.
To evaluate it, use all knot chords with endpoints in
$\Theta=\{0\}\cup\{\theta_c\}$ and the single interval minimum
\begin{equation}\label{na:G-evaluation}
 \min_{\max(0,m-1)\le L\le m}
 [(L+1-m)f(L)+(m-L)\log(L+1)].
\end{equation}
A minimizing law of width less than one can move either endpoint to a
knot: on a piece $f=C+\omega\log x$, an interior stationary chord has
negative second derivative, unless it is constant. Otherwise its width
is one; its upper endpoint is then at least one. This proves the reduction,
including degenerate laws. Shifting and truncating a law shows that $G$
is nondecreasing and $\log G$ is $1/\theta_0$-Lipschitz. The objective in
\eqref{na:G-evaluation} is $2/\theta_0$-Lipschitz in $L$, so a rational
mesh of an interval of length at most one gives a polynomial-work oracle.

\subsection{The physical programs}

For one row and head suppress indices and put $q=.1F+.9d$. There are
three nonconstant ordinary programs:
\begin{equation}\label{na:proxies}
 \begin{split}
 A_p(d)&=\max\{(p+.395494177d)/1.00024010,
                   (p+.4d)/1.00032357,(p+.5d)/1.02631852\},\\
 B_p(q)&=\max\{(p+.9q)/1.19092969,(p+q)/1.23726325\},\\
 J_p(F,d)&=m-E(\max\{d_sp,m-b_s\alpha\},b_s\alpha),\qquad
 m=c_sF+d_s(p+d),
 \end{split}
\end{equation}
where $c_s=.27140539$, $d_s=.72149518$, $b_s=c_s+d_s$.
Only the first branch of $A_p$ is unconditional; all other affine
branches require $\alpha\le4p$ and are omitted otherwise. The stored
functions, with constant baselines subtracted, are respectively
$f(S+h)-p$, $.1f(S)+.9(f(S+h)-p)$, and
$c_sf(S)+d_sf(S+h)-d_sp$. Thus the five affine pieces require only two
programs; their maxima preserve exactly the same means and dual constraints.

For $A,G>0$ the loss in \eqref{na:proxies} is
\begin{equation}\label{na:loss}
 E(A,G)=G\varepsilon(A/G),\qquad
 \varepsilon(r)=r+\frac{\log(1/[r\log(1+1/r)])-1}{\log(1+1/r)},
\end{equation}
with $E(0,G)=G$ and $E(A,0)=0$. Nonpositive scores are ignored.
Include the singleton $p$. Only the home edge also has
\begin{equation}\label{na:home-score}
 \widetilde T_H=\max\{p,\min(p+\alpha v(d/\alpha),F/\beta)\},\qquad
 v(r)=\begin{cases}2^r-1&r\le r_*,\\r-\epsilon_*&r\ge r_*,\end{cases}
\end{equation}
where $\beta=1-e^{-1}$, $\ell=\log2$,
$r_*=\log(1/\ell)/\ell$, and $\epsilon_*=1+r_*-1/\ell$.
HOME stores the union contraction. A singleton stores a constant.
A physical score at most $T_{i0}$ may be dropped: moving its mass to the
private option frees column capacity and does not lower the objective.
In particular there are no physical copies of the private guard.
At $p=0$ the original physical kernel equals the private kernel, so any selected
zero-singleton edge may also be released privately.

The auxiliary NEW label is
\begin{equation}\label{simp:new-score}
 N_{ij}=p_{ij}+(2/5)d_{ij},\qquad
 F_i\le(3/40)p_{ij},\quad\alpha_i\le(21/20)p_{ij}.
\end{equation}
Its column cost is two; ordinary programs have cost one.
Repair may additionally use NEW candidates with $\alpha_i+F_i\le5p_{ij}/3$.
These additional candidates impose no constraint on the auxiliary dual.

\subsection{Why every single program is valid}

The intended outcome payoff of an ordinary label is
$\log\mathcal H_{p,\alpha}(f(S),f(S+h))$; HOME may use the stronger
$\log f(S+h)$, and the private payoff is $\log[\alpha H(f(S)/\alpha)]$.
Section~\ref{simp:completion-section} places all these payoffs under one
pointwise bound for the final rematching. Here we verify their scores
under a last-coin law of raw width at most $\alpha$.

For two nested outcomes there are nonnegative $(r,d,o,v,c)$ with
\[
 p=r+o+c,\quad X_0=d+o,\quad X_1=d+o+v+c,\quad
 V_0=p+d,\quad V_1=p+d+v,\quad g=X_1-X_0=v+c.
\]
Indeed these are the successive raw, union, and overlap increments.
Since $g\le\alpha$, protection can be lowered to $g$. Each affine piece
$(p+k[\chi F+(1-\chi)d])/R$ is certified by
\begin{equation}\label{simp:two-point-test}
 \frac{p+k[d+\chi o+q(v+\chi c)]}{R}
 \le\mathcal H_{p,g}(X_0,V_0)^{1-q}\mathcal H_{p,g}(X_1,V_1)^q
 \quad(0\le q\le1).
\end{equation}
Partition the nonnegative cone at all ties of linear factor forms;
conditional pieces also impose $v+c\le4p$. On each cone, the right
side is a homogeneous concave function, being a weighted geometric mean
of linear forms, with exponents summing to one. Subtracting the linear
left side preserves concavity. It therefore suffices to check extreme
rays, enumerated exactly by intersections of four independent partition
hyperplanes in five dimensions, including boundaries and ties.
On a ray the logarithmic difference is $\log(a+bq)-L-\Delta q$,
where $a,b$ are the numerator coefficients,
$L=\log(R\mathcal H_0)$, and $\Delta=\log(\mathcal H_1/\mathcal H_0)$.
Its maximum occurs at an endpoint or at
$q=[1/\Delta-a/b]_0^1$ when $b\Delta>0$. No grid in $q$ is used.
The four conditional pieces share 60 hyperplanes and 9,326 rays.
For the unconditional contraction piece, decrease raw values to
$X_j=V_j-p$ and protection to the union increment; the analogous
three-dimensional enumeration has 727 rays.
Directed arithmetic verifies every resulting comparison.

For the joint program the finite pointwise certificate is
\begin{equation}\label{na:pointwise-joint}
 \mathcal H_{p,0}(X,V)\ge c_sX+d_sV
 \quad\text{for }\max(X,p)\le V\le X+p.
\end{equation}
Normalize $V=1$. The factor ties partition the triangle
$0\le X,p\le1$, $X+p\ge1$ into cells on which the left side minus
the affine right side is concave. Exact enumeration gives 64 vertices,
and directed checks prove \eqref{na:pointwise-joint} there.
For $Y=A+GZ$, $\Pr(Z=1)=q$, its arithmetic minus geometric mean deficit
is $A+qG-A^{1-q}(A+G)^q$. It decreases with $A$, increases with $G$,
and its maximum over $q$ is \eqref{na:loss}, by differentiation.
The preserved joint mean has width at most $b_s\alpha$ and lower endpoint
at least $\max(d_sp,m-b_s\alpha)$. This proves the joint score.
Also $\varepsilon$ is convex and decreasing, as the supremum of the
convex decreasing deficits $r+q-r^{1-q}(r+1)^q$; hence its perspective
$E$ is jointly convex. Taylor's theorem gives
\begin{equation}\label{simp:loss-tail}
 E(A,G)\le G^2/(8A)\qquad(A>0),
\end{equation}
since the gap's second derivative in $G$ is at most $1/(4A)$.

The private score follows immediately from \eqref{na:private}.
For a HOME law $p+a+bZ$, $b\le\alpha\le p$, and $\ell\le c\le1$,
\[
 \exp\mathbb E\log(p+a+bZ)\ge p+c(a+qb)-\alpha\epsilon(c),\quad
 \epsilon(c)=1+\frac c\ell\left(\log\frac c\ell-1\right).
\]
Removing $a$ and decreasing $p$ can only increase the deficit; its
convex dependence on $b/\alpha$ reduces to the endpoints. At
$p=b=\alpha$ the remaining maximum is
$\max_q(1+cq-2^q)=\epsilon(c)$. Maximizing over $c$ proves $v$.
Clipping only lowers this score. All these bounds are nondecreasing
in the one stored mean, so they remain valid after that mean increases.

The clipping also preserves the HOME score needed by the controller.
With $Q=\max\{0,\max_j(F-d_j-\beta p_j)\}$, set
$d'=(F-Q-\beta p)^+\le d$. If $F/\beta\le p$ the claim is immediate;
otherwise $d'\le F-\beta p$ and $v(r)\le r$ give
$p+\alpha v(d'/\alpha)\le F+e^{-1}p\le F/\beta$.
Thus $\widetilde T_H\ge p+\alpha v(d'/\alpha)$.

\subsection{A finite lower profile, used only in certificates}

The same two-point cone argument in dimension two verifies $G(m)\ge c_0m$.
For lower endpoint at least one all factors equal the raw value, and the
arithmetic-to-geometric ratio decreases with that endpoint, so the last
unbounded interval reduces to its endpoint at one.
For sharper bounds retain 4,001 positive nondecreasing rational targets
$L_i\le G(i/1000)$, $0\le i\le4000$.
Their proof uses the exact reduction \eqref{na:G-evaluation}.
An affine support proposed for the knot chords is checked below $f$ at
every knot; concavity between knots then proves it below the entire chord.
For a full-width law put
$J(L,m)=(L+1-m)f(L)+(m-L)\log(L+1)$. On a piece
$f(L)=C+\omega\log L$,
\[
 L^2(L+1)^2J_{LL}
 =\omega(L+m-1)(L+1)^2-(L+m+2)L^2.
\]
Its four rational Bernstein coefficients certify curvature on an interval.
Concavity or a one-sign enclosure of $J_L$ reduces to endpoint tests;
convexity permits the lower tangent bound
$J(s)-|J_L(s)|(v-u)/2$, $s=(u+v)/2$. Otherwise use a direct interval
bound and bisect. The initial constant piece includes $L=0$.
The complete replay uses 122,316 interval tests, maximum depth 11.

On $[i/1000,(i+1)/1000]$, interpolate $L_{\max(0,i-1)}$ to $L_i$;
after four keep $L_{3999}$ constant. Denote this delayed profile by
$\underline G$. Monotonicity gives $\underline G(m)\le G(m)$ without
assuming concavity of $G$. A certified lower envelope is
\begin{equation}\label{fd:private-lower}
 P(m)=\max\{c_0m,\underline G(m),[m-1/(8(m-1))]^+\ (m>1)\},
\end{equation}
where the last entry is zero for $m\le1$ and follows from
\eqref{simp:loss-tail}. Thus $DP(t/D)\le1$ is a valid relaxation of
$DG(t/D)\le1$. The exact algorithm uses \eqref{na:private}; the finite
profile and all ray enumerations belong solely to its proof.
